\documentclass[aspectratio=169,10pt]{beamer}
\usetheme{default}
\usecolortheme{beaver}
\setbeamertemplate{navigation symbols}{}
\usepackage[utf8]{inputenc}
\usepackage[T1]{fontenc}
\usepackage[french]{babel}
\usepackage{amsmath,amssymb}

\begin{document}

\begin{frame}{Points à aborder}
\vspace{-0.5em}
\small
\begin{enumerate}
  \setlength{\itemsep}{0.2em}
  \item \textbf{Terminologie polytope local / polytope marginal}
  \begin{itemize}
    \item Peut-on l'employer pour des polytopes qui intègrent les contraintes sur la dynamique~?
    \item Faut-il distinguer les variables du polytope local (groupes de variables dont on connaît la dynamique)
          et celles du programme linéaire (degrés de liberté)~?
  \end{itemize}
  \item \textbf{Formulation compacte}
  \begin{itemize}
    \item Comment l'interpréter~?
    \item Faut-il introduire des notations abstraites/explicites pour la forme MILP du problème
          qui consiste à trouver un planning de coût minimum~?
  \end{itemize}
  \item \textbf{Quelles propositions énoncer~?}
  \begin{itemize}
    \item On a l'inclusion des polyèdres.
    \item Mais on ne peut pas comparer les PL, car l'objectif n'est pas le même.
  \end{itemize}
  \item \textbf{Notations} \\
        Les idées à exprimer sont simples. Comment simplifier les notations~?
  \item \textbf{Structure du modèle} \\
        Comment clarifier le passage d'un modèle à trois niveaux (système, substation, composant)
        à une réécriture inspirée des MDP décomposables à deux niveaux (système, sous-systèmes)~?
\end{enumerate}
\end{frame}

\end{document}
